% Folder 138, batch 6 — archivists' pages 101–120.
% Pass: Opus 5.5 (claude-opus-5-5), 2026-09-30 — first pass, unchecked against the pages by a human.
%
% Skipped: none of the twenty sheets is dropped outright. Pages 105 and 109 are
% typed versos (a letter of recommendation dated 11.1.1977; a thesis-defence
% notice of 12.12.1977): only his mathematics on them is transcribed. Pages 112
% and 120 are cover sheets inscribed in his hand; they carry no \page{} and their
% words become section headings.
\documentclass[11pt,a4paper]{article}
\input{../preamble/grothendieck}

\folder{138}
\batch{6}
\pages{101}{120}
\foldertitle{Cartes. Etude arithmétique (1976, 77 ?) : bon de commande (s.d.), notes manuscrites (s.d.).}
\dating{[vers 1976-1978]}
\watermark{Édition de démonstration}

\begin{document}

\page{101}
\note{la page s'ouvre au milieu d'une phrase : l'argument (revêtement galoisien engendré par $f$, groupe $\gamma_{\alpha_0}$, type $\alpha$, entiers $p(s)$, $q(s)$, $r(s)$, $\delta$) vient de la page précédente, hors de ce lot.}

rev.\ étale galoisien \add{de $\mathbb{P}^1_{\Gamma}$} engendré par $f$ \uncertain{sont} de type $\alpha$. Il est évidemment stable par $\gamma_{\alpha_0}$, et correspond à un sous-\struck{\ill{}}\add{schéma} ouvert et fermé $M^{\alpha}_{C_0}$ de $M_{C_0}$, stable par $\gamma_{\alpha_0}$. [Il est vide \uncertain{sauf} si le cardinal $N$ de $\alpha$ est un multiple commun des $p(s)$, $q(s)$, $r(s)$ et de $\delta$ ; \struck{\ill{}} et il est non vide ssi \struck{$\exists$} $G$ \add{de type $\alpha$} \struck{est} engendrable par des gén.\ \struck{$\rho$, $\sigma$} $\rho_s$, $\rho_f$ \add{$\sigma$} satisfaisant
\[
  \rho_s\,\rho_f\,\sigma = 1, \qquad \sigma^2 = 1,
\]
et si \struck{$\exists$} $\exists$ $H$ sous-groupe de $G$, \uncertain{avec} $\bigcap$ conjugués de $H = \{e\}$, et $C_0$ déduit de la carte combinatoire orientée $E = G/H$ \ill{} $\rho_s, \rho_f, \sigma$ \ill{} \ldots]. Mais il est important de réaliser que dans \uncertain{un cas qui arrivât} $N = \delta \mid$ i.e.\ $H = \{e\}$, la classe d'isom.\ des $G$ \emph{munis de} $\rho_s$, $\rho_f$, $\sigma$ (i.e.\ le type de carte combinatoire orientée correspondante) \emph{n'est pas} déterminée par la seule connaissance de $X \xrightarrow{f} \mathbb{P}^1_{\Gamma} \ldots$]\note{l'indice de $\mathbb{P}^1$, ici et dans l'ajout de la première ligne, est lu $\Gamma$ sans certitude.}

\page{102}
\marginal{$\overline{\mathbb{Q}}$ = clôture alg.\ de $\mathbb{Q}$ dans $\mathbb{C}$}
\note{la note marginale est reliée par un trait à un $\overline{\mathbb{Q}}$ écrit au-dessus de « Hom » ; on ne voit pas s'il remplace $\mathbb{C}$ dans $\mathrm{Hom}(R^{\alpha}_{C_0}, \mathbb{C})$.}

\textbf{III)} \uncertain{Considérations} \add{$M^{\alpha}_{C_0}(\overline{\mathbb{Q}}) \simeq$} $M^{\alpha}_{C_0}(\mathbb{C})$ [$= M^{\alpha}_{C_0\,\mathbb{Q}}(\mathbb{C}) \simeq \mathrm{Hom}(R^{\alpha}_{C_0}, \mathbb{C})$, où $R^{\alpha}_{C_0}$ est l'algèbre finie étale sur $\mathbb{Q}$ qui représente $M^{\alpha}_{C_0\,\mathbb{Q}}$]. On sait que c'est aussi l'ens.\ $C^{\alpha}_{C_0}$ des \add{classes d'iso} cartes \add{comb.} orientées qui sont de type $(C_0, \alpha)$, i.e.\ Donc $\Gamma = \mathrm{Aut}(\overline{\mathbb{Q}})$ opère dans $N^{\alpha}_{C_0}$ --- en fait \add{l'} \struck{\ill{}} dans $\Gamma$-ens : \struck{op.} $C^{\alpha}_{C_0}$ s'identifie \ill{} (à iso.\ can.\ près) \ill{} $\mathbb{Q}$-schéma $M^{\alpha}_{C_0\,\mathbb{Q}}$ (ou du $\mathbb{Z}[\frac{1}{\delta!}]$-schéma $(M^{\alpha}_{C_0})_{\mathbb{Z}[\frac{1}{\delta!}]}$ \ldots). Bien entendu, $\Gamma$ opère de façon compatible \add{(i.e.\ commutant)} avec les opérations de $\gamma$ sur $C^{\alpha}_{C_0}$. Rien qu'en regardant l'élément $\tau$ de $\Gamma$, $\tau z = \bar z$, on voit que l'opération de $\Gamma$ sur $C^{\alpha}_{C_0}/\gamma$ ne peut pas être triviale\ldots

On va essayer de déterminer l'opération de $\gamma \times \Gamma$ sur $C^{\alpha}_{C_0}$ dans certains cas élémentaires \uncertain{ou} \ill{} (cartes sphériques \ldots)

\page{103}
\note{page de cas. À gauche de chaque cas, un symbole numérique (ici transcrit en tête de cas) ; à droite de chaque dessin, une colonne où il note le point $\infty$ (« $\infty$ / $\circ$ / multiplicité »). Dans les dessins, $\bullet$ marque les points au-dessus de $0$, $\times$ ceux au-dessus de $1$ ; les chiffres écrits dessous sont les multiplicités. La colonne de droite porte, pour les premiers cas, des égalités $G \simeq \ldots$, $\gamma = \ldots$ lourdement biffées.}

(1) \struck{Carte} combinatoire type \quad [symbole : $1\;1\;1$]

\drawing{un segment, $\bullet$ en $0$ (mult.\ $1$), $\times$ en $1$ (mult.\ $1$) ; $\infty$ de multiplicité $1$.}

$f(z) = z$ \qquad $G \simeq 1$ \struck{$\gamma = 1$} \struck{\ill{}}

(2) Carte combinatoire type \quad [symbole : $2 \begin{smallmatrix}1\\1\end{smallmatrix} 2$]

\drawing{un segment $-1 \times$ --- $\bullet\, 0$ --- $\times\, 1$, multiplicités $1$, $2$, $1$ ; $\infty$ de multiplicité $2$.}

$f(z) = z^2$ \qquad $C_{C_0} = C^{\alpha}_{C_0}$ n'a qu'un élément \struck{\ill{}} \qquad $G \simeq \mathbb{Z}_2$ \struck{$\gamma = \ldots$} \struck{$\gamma = \mathbb{Z}_2$}

(2') Carte comb.\ type \quad [symbole : $\begin{smallmatrix}1\\1\end{smallmatrix} 2\; 2$]

\drawing{un segment $\bullet\,0$ --- $\times\,\frac12$ --- $\bullet\,1$, multiplicités $1$, $1$ ; $\infty$ de multiplicité $2$.}

\[
  \Bigl[\, f(z) = \lambda z(z-1) \qquad
  \begin{cases} f(t) = \lambda t(t-1) = 1 \\ f'(t) = \lambda(2t-1) = 0 \end{cases}
  \quad \text{i.e. } t = \tfrac12,\ \lambda = -4 \,\Bigr]
\]
\[
  f(z) = -4z(z-1)
\]
$C^{\alpha}_{C_0} = C_{C_0}$ n'a qu'un élément. \qquad $G \simeq \mathbb{Z}_2$ \struck{$\gamma = \ldots$}

(\struck{2'')} \quad [symbole : $2\; 2 \begin{smallmatrix}1\\1\end{smallmatrix}$]
\note{un premier dessin (points $0$ et $1$ de multiplicité $2$) et « $f(z) = z^2$ » sont biffés, ainsi que « $G \simeq \mathbb{Z}_2 \simeq \ldots$ ».}

\drawing{une boucle issue de $\bullet\,0$ (mult.\ $2$) et passant par $\times$ en $t = 2$ (mult.\ $2$) ; à l'intérieur, le point $1$ ($\circ$, mult.\ $1$) ; à l'extérieur, $\infty$ ($\circ$, mult.\ $1$).}

\[
  \Bigl[\, f(z) = \lambda \frac{z^2}{z-1}, \qquad f(t) = \lambda\frac{t^2}{t-1} = 1
\]
\[
  \Bigl[\, f'(z) = \lambda\Bigl(\frac{2z(z-1) - z^2}{(z-1)^2}\Bigr) = \lambda\frac{z^2 - 2z}{(z-1)^2} = 0 \,\Bigr]
\]
\[
  f'(t) = 0, \quad t = 2, \quad \lambda = \frac14 \,\Bigr]
\]
\[
  f(z) = \frac14\,\frac{z^2}{z-1}
\]
$C^{\alpha}_{C_0}$ a un seul él.

(3*) \quad [symbole : $3 \begin{smallmatrix}1\\1\\1\end{smallmatrix} 3$]

\drawing{une étoile à trois branches : $\bullet\,0$ au centre (mult.\ $3$), trois $\times$ au-dessus de $1$ (mult.\ $1$) aux extrémités ; $\infty$ de multiplicité $3$.}

$f(z) = z^3$ \qquad $G = \mathbb{Z}_3 \subset \gamma \simeq \mathfrak{S}_3$

$C^{\alpha}_{C_0}$ a deux éléments, $\gamma$ y est transitif, $\Gamma$ opère via la \uncertain{représentation} cyclotomique \uncertain{dans} \ill{}

(3') \quad [symbole : $3 \begin{smallmatrix}1\\2\end{smallmatrix} \begin{smallmatrix}1\\2\end{smallmatrix}$]

\drawing{un segment de $\times$ en $t = -3$ à $\bullet\,0$, puis une boucle issue de $0$ passant par $\times$ en $s = \frac32$ ; à l'intérieur de la boucle, le point $1$ (biffures) ; $\infty$ à l'extérieur.}

\[
  f(z) = \frac{4}{27}\,\frac{z^3}{z-1}
\]
$G \simeq \mathfrak{S}_3$ \qquad $\gamma = \{1\}$ \qquad Card $C^{\alpha}_{C_0} = 1$

\page{104}
\textbf{NB} La figure \uncertain{barycentrique} complète est la réunion de la droite $y\,(= \operatorname{Im} z) = 0$, et de la courbe
\[
  2x(x^2+y^2) + (-3x^2+y^2) = 0
\]
\drawing{l'axe réel, marqué $\times$ en $-3$, $\bullet$ en $0$, $\bullet$ en $1$, $\times$ en $\frac32$, le segment $[0, \frac32]$ hachuré ; la courbe forme une boucle de $0$ à $\frac32$ autour de $1$, et deux branches issues de $0$ qui partent vers le haut et le bas à gauche.}

3'') (dual du précédent)

\drawing{un segment $\bullet\,0$ --- $\times\,\frac13$ --- $\bullet\,1$ --- $\times\,\frac43$ ; $\infty$ à part.}

\[
  f(z) = \frac{27}{4}\,z(z-1)^2
\]
\struck{\ill{}} $G = \mathfrak{S}_3$, \quad $\gamma = \{1\}$ \qquad Card $C^{\alpha}_{C_0} = 1$

\page{105}
\note{ses notes occupent le recto d'une lettre dactylographiée qu'il a biffée ; seules les notes sont transcrites.}

\drawing{deux arbres. (I) : $c \bullet$ --- $\times u$ --- $\bullet a$ --- $\times v$ --- $\bullet b$ --- $\times y$, avec une branche de $a$ vers $\times x$ ; à côté, une flèche courbe et « o.s. ». (II) : $\bullet a$, d'où partent deux branches vers $\times x$ et $\times y$, puis $a$ --- $\times u$ --- $\bullet b$ --- $\times v$ --- $\bullet c$ ; « o.s. ».}

\[
  \begin{matrix} 3 & 2_2 & \\ 2 & 1_2 & 6 \\ 1 & & \end{matrix}
\]
\[
  \sigma_2\sigma_1\sigma_1\sigma_0 = \qquad
  \rho_s\,\rho_f = \sigma, \qquad \rho_f = \rho_s^{-1}\sigma
\]

(I)
\[
  \begin{array}{llll}
  r_1 = (a, ab) & \rho_s r_1 = r_2 & \sigma r_1 = r_4 & \rho_f r_1 = r_5 \\
  r_2 = (a, axa) & \rho_s r_2 = r_3 & \sigma r_4 = r_1 & \rho_f r_5 = r_4 \\
  r_3 = (a, ac) & \rho_s r_3 = r_1 & \sigma r_2 = r_2 & \rho_f r_4 = r_3 \\
  r_4 = (b, ba) & \rho_s r_4 = r_5 & \sigma r_3 = r_6 & \rho_f r_3 = r_6 \\
  r_5 = (b, byb) & \rho_s r_5 = r_4 & \sigma r_6 = r_3 & \rho_f r_6 = r_2 \\
  r_6 = (c, ca) & \rho_s r_6 = r_6 & \sigma r_5 = r_5 & \rho_f r_2 = r_1
  \end{array}
\]
\drawing{sur la ligne $1\;2\;3\;4\;5\;6$ : pour $\rho_s$, une flèche $1 \to 2 \to 3$, une double flèche $4 \leftrightarrow 5$, un point gras en $6$ ; pour $\sigma$, des crochets reliant $1$ et $4$, $3$ et $6$, des points en $2$ et $5$. À droite : $\mathbb{Z}/6\mathbb{Z}$.}

\[
  p = 6 \quad q = 2 \quad r = 6 \qquad N =
\]
\note{« $6$ » et « $2$ » sont écrits en surcharge ; « $N =$ » reste sans valeur.}

(II)
\[
  \begin{array}{llll}
  r_1 = (a, ab) & \rho_s r_1 = r_2 & \sigma r_1 = r_4 & \rho_f r_1 = r_5 \\
  r_2 = (a, axa) & \rho_s r_2 = r_3 & \sigma r_4 = r_1 & \rho_f r_5 = r_6 \\
  r_3 = (a, aya) & \rho_s r_3 = r_1 & \sigma r_2 = r_2 & \rho_f r_6 = r_4 \\
  r_4 = (b, ba) & \rho_s r_4 = r_5 & \sigma r_3 = r_3 & \rho_f r_4 = r_3 \\
  r_5 = (b, bc) & \rho_s r_5 = r_4 & \sigma r_5 = r_6 & \rho_f r_3 = r_2 \\
  r_6 = (c, cb) & \rho_s r_6 = r_6 & \sigma r_6 = r_5 & \rho_f r_2 = r_1
  \end{array}
\]
\drawing{même schéma sur $1\;2\;3\;4\;5\;6$ : $\rho_s$ : $1 \to 2 \to 3$, $4 \leftrightarrow 5$, $6$ fixe ; $\sigma$ : crochet de $1$ à $4$, points en $2$ et $3$, $5 \leftrightarrow 6$.}

\page{106}
\note{le haut de la page est barré de longs traits diagonaux.}

\struck{$\tau_{14}$, $\tau_{43}$ ou $\tau_{46}$ ?} \quad \struck{$\tau_{13}$ --- $\tau_{16}$} \quad \struck{$1 \leq i \leq 5$} \quad \struck{$\rho_s^{i}\sigma\,\rho_s^{j}\sigma \cdots \rho_s^{h}\sigma$}

\struck{$\rho_s^{3} = (4,5)$} \quad \struck{$\rho_s^{4} = (1,2,3)(4)(5)(6)$}

\struck{$(fgh)' = (fg)'h + (fg)h' = f'gh + fg'h + fgh'$}

\[
  a = 0 \quad b = 1 \qquad [\ldots] = \infty
\]
\note{dans les formules, $[\ldots]$ tient la place d'un signe illisible ou biffé illisible (\ill{}), là où l'apparat ne peut entrer dans la formule.}
\uncertain{restent} $c, x, y, u, v$
\[
  f(z) = \lambda\, z^3 (z-1)^2 (z-c)
  \qquad
  \begin{cases}
    0 \text{ zéro triple} \\
    1 \text{ zéro double} \\
    c \text{ zéro simple} \\
    \infty \text{ pôle d'ordre } 6
  \end{cases}
\]
\[
  f'(z) = \lambda z^2(z-1)\bigl[\,\underbrace{3(z-1)(z-c) + 2z(z-c)}_{(5z-3)(z-c)} + z(z-1)\,\bigr]
\]
\[
  = \lambda z^2 (z-1)\bigl[\,6z^2 - (5c+4)z + 3c\,\bigr]
\]
Ses zéros distincts de $0, 1, \infty$ sont les zéros de
\[
  6z^2 - (5c+4)z + 3c = 0
\]
\struck{soient} $u$, $v$. Il faut donc \struck{écrire $f$}
\[
  \begin{cases}
    u, v \text{ distincts entre eux et } \neq 0, 1, c \\
    f(u) = f(v) = 1
  \end{cases}
\]
(racines \ill{} de cette équation \uncertain{?})
\[
  f(z) - 1 = [\ldots]\ \lambda (z-u)^2 (z-v)^2 (z-x)(z-y)
\]

\page{107}
\[
  \frac{1}{\lambda} f(u) = u(\alpha u + \beta)\bigl((\alpha-2)u + (\beta+1)\bigr)(u-c)
\]
\[
  u^2 = \alpha u + \beta, \qquad \alpha = \frac{5c}{6} + \frac23, \quad \beta = -\frac{c}{2}
\]
\note{la première ligne est écrite en surcharge sur une version antérieure (des $u$ biffés) ; deux accolades renvoient aux développements qui suivent.}
\[
  u(\alpha u + \beta) = (\alpha^2+\beta)u + \alpha\beta
\]
\[
  \bigl((\alpha-2)u + (\beta+1)\bigr)(u-c) = (\alpha-2)u^2 + \bigl[(\beta+1) - c(\alpha-2)\bigr]u - (\beta+1)c
\]
\[
  = \bigl(\underbrace{(\alpha-2)\alpha + \beta+1 - c(\alpha-2)}_{(\alpha-2)(\alpha-c) + (\beta+1)}\bigr)u + \bigl((\alpha-2)\beta - (\beta+1)c\bigr)
\]
\[
  \alpha(\alpha^2+\beta)\bigl[(\alpha-2)(\alpha-c)+\beta+1\bigr]\ \struck{\ill{}}
  + \beta(\alpha^2+\beta)\bigl[(\alpha-2)(\alpha-c)+\beta+1\bigr]
\]
\[
  + \alpha\beta\bigl[(\alpha-2)(\alpha-c)+\beta+1\bigr]u
  + (\alpha^2+\beta)\bigl[(\alpha-2)\beta - (\beta+1)c\bigr]u
\]
\[
  + \alpha\beta\bigl[(\alpha-2)\beta - (\beta+1)c\bigr]
\]
\note{la ligne suivante est traversée d'un trait, sans doute une rature.}

\struck{$\bigl[(\alpha-2)(\alpha-c)+\beta+1\bigr]\bigl[(\alpha^2+\beta)(\alpha+1) + \alpha\beta\bigr]$}
\[
  \alpha - 2 = \frac{5c}{6} - \frac43, \qquad \alpha - c = -\frac{c}{6} + \frac23
\]
\[
  (\alpha-2)(\alpha-c) = -\frac{5}{36}c^2 + \frac79 c - \frac89
\]
\[
  \beta + 1 = -\frac{c}{2} + 1 = -\frac{9c}{18} + \frac99
\]
\note{le dernier membre est lu sans certitude ; le numérateur $9c$ est surchargé.}
\[
  (\alpha-2)(\alpha-c) + \beta+1 = -\frac{5}{36}c^2 + \frac{5}{18}c + \frac19
\]
\[
  = -\frac{1}{36}(5c^2 - 10c - 4)
\]
\note{les signes de la dernière parenthèse sont écrits en surcharge.}
\struck{$\alpha^3 + 2\alpha\beta = \frac{(5c+4)^3}{6^3} - \frac{c(5c+4)}{6} = \frac{1}{6^3}\bigl[(5c+4)^3 - 36c(5c+4)\bigr]$}
\note{la ligne précédente est barrée de deux traits ; en dessous, biffé, « $-165c^2$ ».}
\[
  = \alpha(\alpha^2 + 2\beta) = \frac{1}{6^3}(5c+4)\bigl((5c+4)^2 - 36c\bigr)
\]
\[
  = \frac{1}{6^3}(5c+4)(25c^2 + 4c + 16)
\]
\[
  (5c+4)(-29c+4) \qquad\qquad -\frac16\Bigl(\frac56\Bigr)^4 c^5
\]
\note{ces deux dernières expressions sont isolées à droite, sans lien écrit avec le calcul.}

\page{108}
\note{page au crayon.}
\[
  \lambda\bigl(P(c)u + Q(c)\bigr) = 1, \qquad \lambda\bigl(P(c)v + Q(c)\bigr) = 1
\]
\[
  P(c)(u-v) = 0 \qquad \boxed{P(c) = 0} \qquad \lambda = \frac{1}{Q(c)} \qquad Q(c) \neq 0
\]
\drawing{un hexagone, sommets marqués $0$ et $1$, avec des traits qui en partent ; à côté, « $\widetilde{X}$ » et « $\underline{a}$ ».}

\textbf{P} \quad
$\boxed{c \neq 0, 1 \quad Q(c) \neq 0 \quad P(c) = 0 \quad (5c+4)^2 - 72c \neq 0}$
\[
  f(z) = \frac{1}{Q(c)}\, z^3 (z-1)^2 (z-c)
\]
$u, v$ sont les racines de
\[
  z^2 - \frac{5c+4}{6}z + \frac{c}{2} = 0
\]
$x, y$ sont les racines de
\[
  \frac{f(z)-1}{(z-u)^2(z-v)^2} = \frac{f(z)-1}{\bigl(z^2 - \frac{5c+4}{6}z + \frac{c}{2}\bigr)^2}
\]
\drawing{au crayon : une sphère (ellipse) portant un segment de $0$ à $1$ et un point $\infty$ ; un polygone à sept ou huit côtés découpé en secteurs par des rayons issus du centre.}

\page{109}
\note{notes et dessins tracés sur un avis de soutenance dactylographié ; seuls les siens sont transcrits.}
\[
  2\Bigl(\sum p + \sum q + \sum r\Bigr) = 6N
\]
\drawing{à l'encre, un petit arbre inscrit dans une lentille, avec des arêtes vers l'extérieur. Au crayon, sur la moitié inférieure : une carte à faces hachurées, sommets marqués $2$, $4$, $8$, $2$, $4$, $8$, $2$, un segment pointillé ; des courbes emboîtées qui se referment vers un sommet marqué $8$ ; plus bas, un axe allant vers $\infty$ et une figure de cercles concentriques traversée par un segment.}

\page{110}
\[
  f(z) - 1 = \lambda z^3 (z-1)^2 (z-c) - 1 \ \text{divisible par}\ (z^2 - \alpha z - \beta)^2
\]
\[
  \lambda z^3 (z-1)^2 (z-c) - 1 = \lambda (z^2 - \alpha z - \beta)^2 (z^2 + \ell z + m)
\]

$c = 0$ \quad $z^2 - \frac23 z = 0$ \quad d'où $u = 0$, $v = \frac23$
\[
  \frac{1}{\lambda} f = \frac{z^4 (z-1)^2}{z^3 (z-1)^3}
\]
\note{lecture incertaine de cette fraction, dont l'exposant $4$ est en surcharge ; à côté, biffé : « $f = \lambda z^4(z-1)$ ».}
\[
  z^2 - \frac32 z + \frac12 = 0, \qquad (z-1)\Bigl(z - \frac12\Bigr)
\]
\note{un « $6$ » biffé précède $z^2$.}
\struck{$u\,P(c) + Q(c) = 0$} \qquad \struck{$v\,P(c) + Q(c) \neq 0$}
\[
  P(c) = 0 \Longrightarrow c \neq 0, 1
\]
\[
  P(c) = 0 \qquad \underbrace{25c^2 - 32c + 16}_{\delta(c)}
\]
\[
  32(32 - 50) = -32 \cdot 18 = -(8 \cdot 3)^2 \qquad c = \frac{16}{25} \pm \frac{24}{25}\,i
\]
\note{le calcul du discriminant est écrit en surcharge ; la lecture « $32(32-50)$ » est incertaine.}
\[
  \mathbb{Q}[c]_{\delta(c)\,c(c-1)} / P(c)Q(c)
  \qquad
  \mathbb{Q}[c]_{\delta(c)Q(c)}
\]
\[
  P(c) = Q(c) = 0 \Longrightarrow f(u) = f(v) = 0
\]
\begin{itemize}
  \item a) $u = 0 \Longrightarrow c = 0$ absurde \quad ($u, v$ font \uncertain{partie des autres} que $0, 1, c$)
  \item b) $u = 1 \Longrightarrow 6 - (5c+4) + 3c = 0$, i.e.\ $2 - 2c = 0$, i.e.\ $c = 1$ absurde
  \item c) $u = c$ \quad $6c^2 - (5c+4)c + 3c = 0$, \quad $c^2 - c = 0$, \quad $c(1-c) = 0$, \quad $c = 0$ ou $c = 1$ absurde
\end{itemize}

\page{111}
Donc on sait que les zéros de $P(c)$ sont distincts de $0, 1$ et ne sont pas des zéros de $Q$. Mais pourraient-ils être zéros de $\delta(c) = 25c^2 - 32c + 16$ ?

\note{à gauche, le symbole $\begin{smallmatrix}3\\2\\1\end{smallmatrix}$, puis une colonne entourée en surcharge et « $6$ » ; en dessous, « ou » suivi d'un second symbole entouré, biffé.}
\[
  f'(z) = 0 \Longrightarrow f(z) = 0
\]
\note{le dernier signe, un ovale noirci, est lu $0$ sans certitude.}
\[
  f(z) = \lambda\frac{\varphi}{\psi}
\]
$\varphi$, $\psi$ \add{unitaires} premiers entre eux, à zéros indéterminés de multiplicités données, $\deg\varphi > \deg\psi$.
\[
  \frac{f'(z)}{f(z)} = [\ldots]\ \Bigl(\frac{\varphi'}{\varphi} - \frac{\psi'}{\psi}\Bigr)
\]
\[
  \sum \frac{p_i}{z - \alpha_i} - \sum \frac{q_j}{z - \beta_j}
  = \frac{1}{\prod (z-\alpha_i) \prod (z-\beta_j)} \bigl(\qquad\bigr)
\]
\note{la parenthèse reste vide.}

\section{Cartes sphériques en bas degrés}
\note{titre écrit de sa main sur une feuille de couverture (p.~112), qui ne porte rien d'autre.}

\page{113}
\note{dans les tables qui suivent il note le type d'une carte par un triplet ; chaque entrée est une liste d'entiers, écrite en colonne, un indice marquant la répétition ($1_3$ : trois fois $1$). Les colonnes sont rendues ici par $\begin{smallmatrix}a\\b\end{smallmatrix}$. Une accolade reliant deux cartes est transcrite « (accolade) ».}

\drawing{haut de page : une douzaine de petites cartes dessinées à l'encre, certaines entourées d'une courbe fermée portant un nombre ($8$, $9$, $6$) ; des faces marquées $1$, $3$, $4$, $5$ ; des chiffres $3$, $4$, $5$ écrits à côté, en partie surchargés ; un arbre, deux croix à quatre et trois branches ; des esquisses biffées.}

\drawing{un segment à un sommet gras} \quad degré 1
\[
  n = \sum_s q_s = \sum_f p_f
\]
\begin{itemize}
  \item \emph{degré 1} \quad $\varphi(\struck{\ill{}}1) = 2$ \quad \drawing{un segment.}
  \item \emph{degré 2} \quad a) 1 seul sommet d'ordre 2 ; \quad b) 2 sommets d'ordre 1.
  \drawing{pour a) : un segment à sommet au milieu ; une boucle ; (biffé) ; pour b) : un segment à deux sommets.}
  \item \emph{degré 3} \quad a) 1 sommet d'ordre 3 ; \quad b) 2 sommets, d'ordre 1 et 2.
  \drawing{un tripode ; une boucle avec une queue ; une esquisse biffée ; un segment à deux sommets prolongé.}
  \item \emph{degré 4} \quad a) 1 sommet d'ordre 4 ; \quad b) 1 sommet d'ordre 3, 1 d'ordre 1.
  \drawing{pour a) : une croix ; une boucle traversée d'un segment ; deux boucles accolées ; une boucle avec un segment intérieur (ces deux dernières séparées par un trait courbe). Pour b) : un segment suivi d'une fourche ; des esquisses biffées ; « rigide ».}
\end{itemize}

\page{114}
\note{recensement par degré ; sous chaque carte, son groupe d'automorphismes (souligné deux fois quand il l'a souligné), puis son type.}

(1) \quad \drawing{un segment.} \quad $\mathbb{Z}_2$ \quad $(1, 1, 1)$

(2) \quad \drawing{un segment à sommet au milieu ; un segment à deux sommets ; une boucle.} \quad $\mathbb{Z}_2 \times \mathbb{Z}_2$, \quad $\mathbb{Z}_2 \times \mathbb{Z}_2$, \quad $\mathbb{Z}_2 \times \mathbb{Z}_2$ (les deux dernières sous une accolade)
\[
  \bigl(2, \begin{smallmatrix}1\\1\end{smallmatrix}, 2\bigr) \quad
  \bigl(\begin{smallmatrix}1\\1\end{smallmatrix}, 2, 2\bigr) \quad
  \bigl(2, 2, \begin{smallmatrix}1\\1\end{smallmatrix}\bigr)
\]

(3) \quad \drawing{un tripode ; une boucle avec une queue ; un segment à deux sommets prolongé.} \quad $\mathfrak{S}_3$, \quad $\mathbb{Z}_2$, \quad $\mathbb{Z}_2$ (les deux dernières sous une accolade)
\[
  \bigl(3, \begin{smallmatrix}1\\1\\1\end{smallmatrix}, 3\bigr) \quad
  \bigl(3, \begin{smallmatrix}2\\1\end{smallmatrix}, \begin{smallmatrix}2\\1\end{smallmatrix}\bigr) \quad
  \bigl(\begin{smallmatrix}2\\1\end{smallmatrix}, \begin{smallmatrix}2\\1\end{smallmatrix}, 3\bigr)
\]

(4)
\drawing{un segment suivi d'une fourche ; une boucle avec deux arêtes pendantes au même sommet (accolade).} \quad $\mathbb{Z}_2$, $\mathbb{Z}_2$
\[
  \bigl(\begin{smallmatrix}3\\1\end{smallmatrix}, \begin{smallmatrix}2\\1\\1\end{smallmatrix}, 4\bigr) \quad
  \bigl(4, \begin{smallmatrix}2\\1\\1\end{smallmatrix}, \begin{smallmatrix}3\\1\end{smallmatrix}\bigr)
\]
\drawing{un segment à deux sommets intérieurs ; une boucle traversée d'un segment (accolade).} \quad $\mathbb{Z}_2 \times \mathbb{Z}_2$, $\mathbb{Z}_2 \times \mathbb{Z}_2$
\[
  \bigl(\begin{smallmatrix}2\\2\end{smallmatrix}, \begin{smallmatrix}2\\1\\1\end{smallmatrix}, 4\bigr) \quad
  \bigl(4, \begin{smallmatrix}2\\1\\1\end{smallmatrix}, \begin{smallmatrix}2\\2\end{smallmatrix}\bigr)
\]
\drawing{une croix.} \quad $\underline{\underline{\mathbb{D}_4}}$ \quad $(4, \begin{smallmatrix}1\\1\\1\\1\end{smallmatrix}, 4)$

\drawing{un segment à trois sommets ; deux boucles accolées en un sommet (accolade).} \quad $\mathbb{Z}_2 \times \mathbb{Z}_2$, $\mathbb{Z}_2 \times \mathbb{Z}_2$
\[
  \bigl(\begin{smallmatrix}2\\1\\1\end{smallmatrix}, \begin{smallmatrix}2\\2\end{smallmatrix}, 4\bigr) \quad
  \bigl(4, \begin{smallmatrix}2\\2\end{smallmatrix}, \begin{smallmatrix}2\\1\\1\end{smallmatrix}\bigr)
\]
\drawing{un cercle portant deux sommets.} \quad $\underline{\underline{\mathbb{D}_4}}$ \quad $\bigl(\begin{smallmatrix}2\\2\end{smallmatrix}, \begin{smallmatrix}2\\2\end{smallmatrix}, \begin{smallmatrix}2\\2\end{smallmatrix}\bigr)$

\drawing{plus pâle : un cercle avec une queue terminée par un sommet.} \quad $\mathbb{Z}_2$ \quad $\bigl(\begin{smallmatrix}3\\1\end{smallmatrix}, \begin{smallmatrix}2\\2\end{smallmatrix}, \begin{smallmatrix}1\\3\end{smallmatrix}\bigr)$
\note{les chiffres de ce dernier type sont peu sûrs.}

(5)
\drawing{une étoile à cinq branches.} \quad $\underline{\underline{\mathbb{D}_5}}$ \quad \struck{\ill{}}, $(5, 1_5, 5)$

\drawing{une croix avec un sommet à l'extrémité d'une branche ; une boucle avec trois arêtes pendantes au même sommet (accolade).} \quad $\mathbb{Z}_2$, $\mathbb{Z}_2$
\[
  \bigl(\begin{smallmatrix}4\\1\end{smallmatrix}, \begin{smallmatrix}2\\1_3\end{smallmatrix}, 5\bigr) \quad
  \bigl(5, \begin{smallmatrix}2\\1_3\end{smallmatrix}, \begin{smallmatrix}4\\1\end{smallmatrix}\bigr)
\]
\drawing{une boucle traversée par un segment, avec deux arêtes pendantes ; un segment suivi d'une fourche (accolade).} \quad $\mathbb{Z}_2$, $\mathbb{Z}_2$
\[
  \bigl(5, \begin{smallmatrix}2\\1_3\end{smallmatrix}, \begin{smallmatrix}3\\2\end{smallmatrix}\bigr) \quad
  \bigl(\begin{smallmatrix}3\\2\end{smallmatrix}, \begin{smallmatrix}2\\1_3\end{smallmatrix}, 5\bigr)
\]
\drawing{deux boucles accolées avec une arête intérieure ; un segment à trois sommets prolongé (accolade).} \quad $\mathbb{Z}_2$, $\mathbb{Z}_2$
\[
  \bigl(5, \begin{smallmatrix}2_2\\1\end{smallmatrix}, \begin{smallmatrix}2_2\\1\end{smallmatrix}\bigr) \quad
  \bigl(\begin{smallmatrix}2_2\\1\end{smallmatrix}, \begin{smallmatrix}2_2\\1\end{smallmatrix}, 5\bigr)
\]

\page{115}
\note{suite du degré 5.}

\drawing{deux boucles accolées avec une arête pendante vers le haut ; un tripode à deux sommets terminaux (accolade).} \quad $\mathbb{Z}_2$, $\mathbb{Z}_2$
\[
  \bigl(5, \begin{smallmatrix}2_2\\1\end{smallmatrix}, \begin{smallmatrix}3\\1_2\end{smallmatrix}\bigr) \quad
  \bigl(\begin{smallmatrix}3\\1_2\end{smallmatrix}, \begin{smallmatrix}2_2\\1\end{smallmatrix}, 5\bigr)
\]
\drawing{une boucle traversée par un segment terminé par un sommet ; une boucle avec une queue à deux sommets (accolade).} \quad $\mathbb{Z}_2$, $\mathbb{Z}_2$
\[
  \bigl(\begin{smallmatrix}4\\1\end{smallmatrix}, \begin{smallmatrix}2_2\\1\end{smallmatrix}, \begin{smallmatrix}3\\2\end{smallmatrix}\bigr) \quad
  \bigl(\begin{smallmatrix}3\\2\end{smallmatrix}, \begin{smallmatrix}2_2\\1\end{smallmatrix}, \begin{smallmatrix}4\\1\end{smallmatrix}\bigr)
\]
\drawing{un cercle portant deux sommets, dont l'un a une arête vers l'extérieur.} \quad $\mathbb{Z}_2$ ! \quad $\bigl(\begin{smallmatrix}3\\2\end{smallmatrix}, \begin{smallmatrix}2_2\\1\end{smallmatrix}, \begin{smallmatrix}3\\2\end{smallmatrix}\bigr)$

\drawing{un cercle avec, au même sommet, une queue terminée par un sommet et une arête libre.} \quad $1$ \quad $\bigl(\begin{smallmatrix}4\\1\end{smallmatrix}, \begin{smallmatrix}2_2\\1\end{smallmatrix}, \begin{smallmatrix}4\\1\end{smallmatrix}\bigr)$

\page{116}
\note{types en degré 6, sans dessins. Les doubles barres obliques et les biffures sont les siennes ; plusieurs chiffres sont écrits en surcharge.}
\[
  (6, 1_6, 6) \;/\!/\; \bigl(6, \begin{smallmatrix}2\\1_4\end{smallmatrix}, \begin{smallmatrix}5\\1\end{smallmatrix}\bigr) \quad
  \bigl(6, \begin{smallmatrix}2\\1_4\end{smallmatrix}, \begin{smallmatrix}4\\2\end{smallmatrix}\bigr) \quad
  \bigl(6, \begin{smallmatrix}2\\1_4\end{smallmatrix}, \begin{smallmatrix}3\\3\end{smallmatrix}\bigr) \;/\!/\; [\ldots]
\]
\[
  \bigl(6, \begin{smallmatrix}2_2\\1_2\end{smallmatrix}, \begin{smallmatrix}4\\1_2\end{smallmatrix}\bigr) \quad
  \bigl(6, \begin{smallmatrix}2_2\\1_2\end{smallmatrix}, \begin{smallmatrix}3\\2\\1\end{smallmatrix}\bigr) \quad
  \bigl(6, \begin{smallmatrix}2_2\\1_2\end{smallmatrix}, 2_3\bigr) \quad
  \bigl(6, \begin{smallmatrix}2_2\\1_2\end{smallmatrix}, \begin{smallmatrix}3\\2\\1\end{smallmatrix}\bigr)
\]
\note{une accolade relie le deuxième et le dernier type de cette ligne, qui sont identiques, avec un point d'exclamation.}
\[
  \bigl(6, 2_3, \begin{smallmatrix}3\\1_3\end{smallmatrix}\bigr) \quad
  \bigl(6, 2_3, \begin{smallmatrix}2_2\\1_2\end{smallmatrix}\bigr)
\]
\[
  \bigl(\begin{smallmatrix}5\\1\end{smallmatrix}, \begin{smallmatrix}2\\1_4\end{smallmatrix}, 6\bigr) \quad
  \bigl(\begin{smallmatrix}5\\1\end{smallmatrix}, \begin{smallmatrix}2_2\\1_2\end{smallmatrix}, \begin{smallmatrix}5\\1\end{smallmatrix}\bigr) \quad
  \bigl(\begin{smallmatrix}5\\1\end{smallmatrix}, \begin{smallmatrix}2_2\\1_2\end{smallmatrix}, \begin{smallmatrix}4\\2\end{smallmatrix}\bigr) \quad
  \bigl(\begin{smallmatrix}5\\1\end{smallmatrix}, \begin{smallmatrix}2_2\\1_2\end{smallmatrix}, \begin{smallmatrix}3\\3\end{smallmatrix}\bigr)
\]
\[
  \bigl(\begin{smallmatrix}5\\1\end{smallmatrix}, 2_3, \begin{smallmatrix}4\\1_2\end{smallmatrix}\bigr) \quad
  \bigl(\begin{smallmatrix}5\\1\end{smallmatrix}, 2_3, \begin{smallmatrix}3\\2\\1\end{smallmatrix}\bigr)
\]
\[
  \bigl(\begin{smallmatrix}4\\2\end{smallmatrix}, \begin{smallmatrix}2\\1_4\end{smallmatrix}, 6\bigr) \quad
  \bigl(\begin{smallmatrix}4\\2\end{smallmatrix}, \begin{smallmatrix}2_2\\1_2\end{smallmatrix}, \begin{smallmatrix}5\\1\end{smallmatrix}\bigr) \quad
  \bigl(\begin{smallmatrix}4\\2\end{smallmatrix}, \begin{smallmatrix}2_2\\1_2\end{smallmatrix}, \begin{smallmatrix}4\\2\end{smallmatrix}\bigr)
\]
\[
  \bigl(\begin{smallmatrix}4\\1_2\end{smallmatrix}, \begin{smallmatrix}2_2\\1_2\end{smallmatrix}, 6\bigr) \quad
  \bigl(\begin{smallmatrix}4\\1_2\end{smallmatrix}, 2_3, \begin{smallmatrix}5\\1\end{smallmatrix}\bigr)
\]
\note{dans les deux lignes suivantes, la première entrée est écrite $3_2$ en surcharge sur une colonne antérieure (sans doute $\begin{smallmatrix}3\\2\\1\end{smallmatrix}$).}
\[
  \bigl(3_2, 2_3, \begin{smallmatrix}4\\1_2\end{smallmatrix}\bigr) \quad
  \bigl(3_2, \begin{smallmatrix}2_2\\1_2\end{smallmatrix}, \begin{smallmatrix}5\\1\end{smallmatrix}\bigr) \quad
  \bigl(3_2, \begin{smallmatrix}2\\1_4\end{smallmatrix}, 6\bigr)
\]
\[
  \bigl(3_2, 2_3, 2_3\bigr) \quad
  \bigl(3_2, \begin{smallmatrix}2_2\\1_2\end{smallmatrix}, \begin{smallmatrix}4\\2\end{smallmatrix}\bigr) \quad
  \bigl(3_2, \begin{smallmatrix}2_2\\1_2\end{smallmatrix}, \begin{smallmatrix}3\\3\end{smallmatrix}\bigr)
\]
\note{d'une autre encre, plus claire, cinq types numérotés, qui renvoient aux cartes marquées des mêmes lettres p.~118.}
\[
  (\alpha)\ \bigl(\begin{smallmatrix}4\\2\end{smallmatrix}, \begin{smallmatrix}2_2\\2_1\end{smallmatrix}, \begin{smallmatrix}3_2\\3\end{smallmatrix}\bigr) \qquad
  (\beta)\ \bigl(\begin{smallmatrix}4\\2\end{smallmatrix}, \begin{smallmatrix}[\ldots]\\2_3\end{smallmatrix}, \begin{smallmatrix}3\\2\\1\end{smallmatrix}\bigr) \qquad
  (\gamma)\ \bigl(\begin{smallmatrix}4\\2\end{smallmatrix}, \begin{smallmatrix}2_2\\2_1\end{smallmatrix}, \begin{smallmatrix}4\\2\end{smallmatrix}\bigr)
\]
\[
  (\delta)\ \bigl(\begin{smallmatrix}4\\2_1\end{smallmatrix}, \begin{smallmatrix}2_2\\2_1\end{smallmatrix}, 6\bigr) \qquad
  (\eta)\ \bigl(\begin{smallmatrix}4\\2_1\end{smallmatrix}, 2_3, \begin{smallmatrix}3_2\\ {[\ldots]}\end{smallmatrix}\bigr)
\]
\note{les indices « $2_1$ » sont lus tels qu'écrits ; plusieurs sont en surcharge.}

\page{117}
\note{suite du recensement de la p.~113 (degré 4, puis degré 5), avec les dessins à droite.}

c) 2 sommets d'ordre 2 \drawing{un cercle portant deux sommets ; un segment à deux sommets prolongé des deux côtés.}

d) 3 sommets, dont deux d'ordre 1, 1 d'ordre 2 \drawing{un segment à trois sommets ; séparé par un trait courbe, un tripode dans une boucle, marqué $\times$.}

\emph{degré 5} \struck{\ill{}}
\begin{itemize}
  \item a) 1 sommet d'ordre 5 \drawing{une étoile à cinq branches ; une boucle à trois arêtes pendantes ; des boucles accolées ; plusieurs sont marquées $\times$.}
  \item b) 2 sommets, d'ordre 4 et 1 \drawing{un segment suivi d'un tripode ; un segment suivi d'une boucle avec arête, encadré « rigide » ; une boucle avec une arête intérieure et une queue, marquée $\times$.}
  \item c) 2 sommets, d'ordre 3 et 2 \drawing{un segment suivi d'une fourche, marqué $\times$ ; un segment suivi d'une boucle, marqué $\times$ ; un cercle portant deux sommets avec une arête vers l'extérieur.}
  \item d) 3 sommets, dont 1 d'ordre 3, deux d'ordre 1 \drawing{un tripode à deux sommets terminaux.}
  \item e) 3 sommets, dont 2 d'ordre 2, 1 d'ordre 1 \drawing{une esquisse biffée ; un segment à trois sommets prolongé, marqué $\times$.}
  \item f) 4 sommets, dont trois d'ordre 1, 1 d'ordre 2 \quad impossible \ldots
\end{itemize}

\page{118}
\note{degré 6. Sous chaque carte, son groupe d'automorphismes ; un chiffre entouré ($1$) marque une carte sans automorphisme non trivial. Des traits courbes au crayon regroupent certaines cartes ; les petites cartes marquées $(\alpha)$, $(\beta)$, $(\gamma)$, $(\delta)$, $(\eta)$ renvoient aux types de la p.~116.}

\emph{degré 6} \quad a) 1 sommet d'ordre 6
\drawing{une étoile à six branches, $\underline{\underline{\mathbb{D}_6}}$ ; séparée par une double barre, une étoile à quatre branches avec une boucle, $\mathbb{Z}_2$ ; un tripode avec une boucle, $\mathbb{Z}_2$ ; une boucle traversée à trois arêtes pendantes, $\mathbb{Z}_2 \times \mathbb{Z}_2$ ; deux boucles accolées à deux arêtes, $\mathbb{Z}_2$ ; une boucle contenant un point, avec deux arêtes intérieures, $\mathbb{Z}_2$ ; deux boucles accolées traversées d'un segment, $\mathbb{Z}_2 \times \mathbb{Z}_2$ ; deux boucles et une arête, $1$ ; trois boucles, $\mathfrak{S}_3$ ; une boucle dans une boucle, $\mathbb{Z}_2$ ; au crayon, dans un enclos, deux boucles traversées d'un segment, $\mathbb{Z}_2 \times \mathbb{Z}_2$, avec le type (d'une autre encre) $\bigl(6, \begin{smallmatrix}2_2\\1_2\end{smallmatrix}, \begin{smallmatrix}4\\1_2\end{smallmatrix}\bigr)$.}

b) 1 sommet d'ordre 5, 1 d'ordre 1
\drawing{un segment suivi d'une étoile à quatre branches, $\mathbb{Z}_2$ ; un segment suivi d'une étoile avec boucle, $1$ (un $\mathbb{Z}_2$ biffé) ; un segment suivi d'une boucle avec arêtes, $1$ ; un segment suivi d'une boucle à deux arêtes intérieures, $\mathbb{Z}_2$ ; un segment suivi de deux boucles, $\mathbb{Z}_2$ ; deux boucles et une arête terminée par un sommet, $1$ ; au crayon, dans un enclos, une boucle avec segment, $\mathbb{Z}_2$.}

c) 1 sommet d'ordre 4, 1 d'ordre 2
\drawing{à gauche, petites cartes d'une autre encre : $(\alpha)$ une boucle traversée d'un segment, $\mathbb{Z}_2$ ; $(\gamma)$ une petite croix dans un cercle, $\mathbb{Z}_2$ ; $(\beta)$ deux boucles accolées, $\mathbb{Z}_2$. Puis : un segment à deux sommets suivi d'une fourche, $\mathbb{Z}_2$ ; un segment suivi d'une boucle avec arête, $1$ ; un segment suivi d'une boucle, $\mathbb{Z}_2$.}

d) 1 sommet d'ordre 4, deux d'ordre 1
\drawing{une croix à deux sommets terminaux, $\mathbb{Z}_2$ ; une boucle avec deux arêtes à sommets terminaux, $\mathbb{Z}_2$ ; $(\delta)$ une croix à deux sommets opposés, $\mathbb{Z}_2 \times \mathbb{Z}_2$ ; $(\eta)$ un cercle avec un rayon prolongé, $\mathbb{Z}_2$.}

e) 2 sommets d'ordre 3
\drawing{une esquisse biffée ($\mathbb{Z}$) ; deux cercles reliés par un segment, $\mathbb{Z}_2 \times \mathbb{Z}_2$ ; un cercle avec une queue terminée par une fourche, $\mathbb{Z}_2$ ; deux fourches reliées par un segment, $\mathbb{Z}_2 \times \mathbb{Z}_2$. En bas de page : une ellipse partagée par un diamètre, $\mathbb{Z}_2 \cdot \mathfrak{S}_3$ ; une esquisse biffée ; deux ellipses traversées d'arêtes, $\mathbb{Z}_2 \times \mathbb{Z}_2$ et $\mathbb{Z}_2 \times \mathbb{Z}_2$.}

\page{119}
f) 1 sommet d'ordre 3, 1 d'ordre 2, 1 d'ordre 1
\drawing{un tripode avec un sommet sur une branche, $1$ ; une boucle avec une queue, marquée d'une croix ; des esquisses encadrées de traits et en partie biffées, $1$, $\mathbb{Z}_2$ ; un tripode suivi d'un segment à sommet, marqué $*$, $\mathbb{Z}_2$ ; un cercle avec une queue à deux sommets, $\mathbb{Z}_2$.}
\[
  \bigl(\begin{smallmatrix}3\\2\\1\end{smallmatrix}, \begin{smallmatrix}2_2\\1_2\end{smallmatrix}, 6\bigr) \quad
  \bigl(\begin{smallmatrix}3\\2\\1\end{smallmatrix}, \begin{smallmatrix}2_2\\1_2\end{smallmatrix}, [\ldots]\bigr) \quad
  \bigl(\begin{smallmatrix}3\\2\\1\end{smallmatrix}, 2_3, \begin{smallmatrix}5\\1\end{smallmatrix}\bigr) \quad
  \bigl(\begin{smallmatrix}3\\2\\1\end{smallmatrix}, 2_3, \begin{smallmatrix}4\\2\end{smallmatrix}\bigr)
\]
\note{ces quatre types sont écrits en colonne à droite, en partie sous une grosse tache d'encre ; une accolade regroupe les deux derniers.}

\struck{g} g) 1 sommet d'ordre 3, trois d'ordre 1
\drawing{un tripode à trois sommets terminaux, $\underline{\underline{\mathbb{D}_3}}$ ; une esquisse biffée.} \quad $\bigl(\begin{smallmatrix}3\\1_3\end{smallmatrix}, 2_3, 6\bigr)$

h) 3 sommets d'ordre 2
\drawing{un cercle portant trois sommets, $\underline{\mathbb{Z}_2 \cdot \mathbb{D}_3}$ (un premier symbole biffé) ; un segment à trois sommets prolongé des deux côtés, $\mathbb{Z}_2 \times \mathbb{Z}_2$.}
\[
  (2_3, 2_3, 3_2) \qquad \bigl(2_3, \begin{smallmatrix}2_2\\1_2\end{smallmatrix}, 6\bigr)
\]
\note{la première entrée du second type est écrite en surcharge.}

i) 2 sommets d'ordre 2, 2 d'ordre 1
\drawing{un segment à quatre sommets, $\mathbb{Z}_2 \times \mathbb{Z}_2$.} \quad $\bigl(\begin{smallmatrix}2_2\\1_2\end{smallmatrix}, 2_3, 6\bigr)$

j) 1 sommet d'ordre 2, 4 d'ordre 1 \quad impossible

\section{« Opérations » sphériques sur les cartes (axiomatique provisoire)}
\note{titre écrit de sa main sur une feuille de couverture (p.~120), qui ne porte rien d'autre ; la suite qu'il annonce commence après ce lot.}

\end{document}
